\documentclass[11pt]{article}
\usepackage[margin=1in]{geometry}
\usepackage{amsmath,amssymb,amsthm,mathtools}
\usepackage{booktabs}
\usepackage{xcolor}
\usepackage[colorlinks=true,linkcolor=blue!55!black,citecolor=blue!55!black,urlcolor=blue!55!black]{hyperref}
\usepackage{microtype}

\newtheorem{theorem}{Theorem}[section]
\newtheorem{proposition}[theorem]{Proposition}
\newtheorem{lemma}[theorem]{Lemma}
\newtheorem{corollary}[theorem]{Corollary}
\newtheorem{conjecture}[theorem]{Conjecture}
\theoremstyle{remark}
\newtheorem{remark}[theorem]{Remark}
\newcommand{\eps}{\{-1,1\}}
\newcommand{\E}{\mathbb E}
\newcommand{\Osc}{\mathrm{Osc}}

\title{\textbf{The minimax quadratic sign problem:\\
exact values, oscillation superadditivity,\\ and the two-scale barrier}\\[4pt]
\large A continuation of the study of MathOverflow question 413935}
\author{mo-413935-ai-attempt repository, second research note}
\date{1 August 2026}

\begin{document}
\maketitle

\begin{abstract}
For
\[
 F(n)=\min_{a_{ij}\in\eps}\max_{x\in\eps^n}
 \Bigl|\sum_{1\le i<j\le n}a_{ij}x_ix_j\Bigr|,
\]
the question whether $F(n)/n^{3/2}$ converges remains open. This note
extends the July 2026 note in this repository. New content: the exact
values $F(2),\dots,F(10)=1,3,4,4,5,9,10,12,13$; a proof that the
Gaussian-rounding method of the July note, fully optimized, computes
exactly the oscillation functional
$\Osc_{1/2}(n)=\tfrac12\min_A(\max_xS_A-\min_xS_A)$, which lies in
$[F(n)/2,\,F(n)]$ and is superadditive; a merge theorem
$F(n+m)\le F(n)+F(m)+(1+o(1))\,n\sqrt m$ with a matching block floor of
order $n\sqrt m$; a resulting continuity theory for
$G(n)=F(n)/n^{3/2}$, including the statement that the limit exists if
and only if it exists along conference orders; the observation that
symmetric conference matrices attain their spectral ceiling exactly at
orders $q+1$ with $q$ an even prime power ($q=9,25$ verified), so no
uniform conference deficit exists; and an explicit bounded divergent
profile satisfying every inequality proved here and in the July note,
which shows that all such local inequalities together cannot decide
convergence. The sharpest remaining target is a one-scale-to-all-scales
amplification inequality. The question stays open.
\end{abstract}

\section{The problem and what is new}

Throughout, $A$ is a symmetric matrix with zero diagonal and
off-diagonal entries $a_{ij}\in\eps$, called a signing of $K_n$. Write
\[
 S_A(x)=\sum_{i<j}a_{ij}x_ix_j,\qquad
 M(A)=\max_{x\in\eps^n}|S_A(x)|,\qquad
 F(n)=\min_AM(A),\qquad
 G(n)=\frac{F(n)}{n^{3/2}}.
\]
MathOverflow question 413935, posed by Paata Ivanisvili in January
2022, asks whether $\lim_nG(n)$ exists. The July 2026 note in this
repository proves
\[
 \frac1\pi\le\liminf G\le\limsup G\le\frac12
\]
and the covering-radius identity $F(n)=\binom n2-2\rho(D_n)$. Those
proofs were re-audited for the present note and stand; see
\texttt{AUDIT.md}.

New results proved below.

\begin{enumerate}
\item Exact values $F(2),\dots,F(10)$ and the exact one-sided and
oscillation optima for $n\le10$, none previously recorded
(Section~\ref{sec:exact}).
\item An even-odd inequality: for every signing, every vertex
bipartition, and every pair $(x,y)$,
$M\ge|S_P(x)+S_R(y)|+|x^{\mathsf T}Cy|$, giving a floor of order
$n\sqrt m$ for every two-block structure (Section~\ref{sec:evenodd}).
\item A merge theorem $F(n+m)\le F(n)+F(m)+\sqrt{nm\,q(n)}$ with
$q(n)=(1+o(1))n$, hence a modulus of continuity for $G$ on
multiplicative windows, the interval structure of the limit set, and
the reduction of the limit question to conference orders
(Section~\ref{sec:merge}).
\item The oscillation functional $\Osc(n)$ is superadditive, satisfies
$F\le\Osc\le2F$, and equals the exact value of the July note's
Gaussian-rounding method when that method is optimized over all
correlation matrices; the constant $1/\pi$ is the value of the linear
slice, and no polynomial slice can beat it on flat-spectrum signings
(Section~\ref{sec:osc}).
\item A fixed-temperature equivalence: for every fixed $\beta>0$ the
free-energy version of the problem differs from $G(n)$ by at most
$\log2/(\beta\sqrt n)$, so the zero-temperature limit is harmless and
the entire difficulty is the $n\to\infty$ interpolation
(Section~\ref{sec:beta}).
\item Relaxations collapse: couplings in $[-1,1]^{\binom n2}$ give
optimum $0$; couplings on the Euclidean sphere of radius
$\sqrt{\binom n2}$ give optimum exactly $\sqrt{\binom n2}=\Theta(n)$;
near-optimal signings converge to the zero graphon
(Section~\ref{sec:relax}).
\item Symmetric conference matrices attain their spectral ceiling
exactly at orders $10$ and $26$, and in general at any order $q+1$,
$q$ an even power of an odd prime, whose switching class contains a
suitable regular graph. The July note's ``conference matrices do not
attain their spectral ceiling'' is therefore an artifact of
$\sqrt q\notin\mathbb Z$ at the orders then checked
(Section~\ref{sec:attain}).
\item A consistency barrier: an explicit bounded divergent profile
satisfies every inequality about the values $F(n)$ proved here and in
the July note. Local estimates alone, however sharpened, cannot settle
the question (Section~\ref{sec:barrier}).
\end{enumerate}

Section~\ref{sec:programs} organizes three proof programs with their
exact targets and failure points, and Section~\ref{sec:wall} states
the sharpest remaining wall and the single recommended next lemma.
All finite claims are reproducible from
\texttt{verification/}; see the appendix.

\section{Exact values and structural data}\label{sec:exact}

\begin{proposition}[Parity]\label{prop:parity}
Every value $S_A(x)$, hence $M(A)$ and $F(n)$, is congruent to
$\binom n2$ modulo $2$.
\end{proposition}
\begin{proof}
$S_A(x)$ is a sum of $\binom n2$ odd terms.
\end{proof}

\begin{theorem}[Exact values]\label{thm:exact}
By exhaustive enumeration of switching classes (representatives with
all edges at one vertex positive; the maximum is a switching-class
invariant),
\[
\begin{array}{c|ccccccccc}
n&2&3&4&5&6&7&8&9&10\\\hline
F(n)&1&3&4&4&5&9&10&12&13\\
G(n)&0.354&0.577&0.500&0.358&0.340&0.486&0.442&0.444&0.411
\end{array}
\]
Witness signings are recorded in
\texttt{verification/data/results\_2026\_08.json} and reverified by an
independent script.
\end{theorem}

Remarks on the data.

\begin{enumerate}
\item $G$ is not monotone in either direction and oscillates through
$n=10$. The smallest ratios in the exact range occur at $n=6$ and
$n=5$, the conference two-graph and its one-point deletion; among
$7\le n\le10$ the smallest is the conference order $n=10$.
\item At $n=5,6$ the optimal switching classes are exactly the
conference-derived two-graphs: the Seidel spectra of the optimizers
are $\{\pm\sqrt5^{(2)},0\}$ and $\{\pm\sqrt5^{(3)}\}$.
\item At $n=10$ the optimum $F(10)=13$ is strictly better than the
conference class, whose maximum is $15$ (Section~\ref{sec:attain}),
and the optimal spectrum $\{\pm\sqrt{17},\pm3.387^{(2)},\pm1.590^{(2)}\}$
is far from flat. Spectral flatness and minimax optimality diverge
already at $n=10$.
\item One-sided optimum. Enumeration gives
$\min_A\max_xS_A(x)=\lfloor n/2\rfloor$ for $3\le n\le10$, attained
uniquely by the all-negative class. This matches the theorem of
Petersdorf \cite{petersdorf} that the maximum frustration index of a
signing of $K_n$ is $\lfloor(n-1)^2/4\rfloor$, attained only by the
all-negative class: via the cut-code dictionary the two statements are
equivalent. The one-sided problem is therefore $\Theta(n)$, while the
two-sided problem is $\Theta(n^{3/2})$: the absolute value is the
entire content of the $n^{3/2}$ scale. In coding terms, the cut code
has covering radius $\binom n2/2-\Theta(n)$, while appending the
all-ones word (the code $D_n$ of the July note is the cut code
extended by complementation, a linear code) drops the covering radius
to $\binom n2/2-\Theta(n^{3/2})$.
\item Certified upper bounds beyond the exact range, from principal
submatrices of Paley conference matrices, from the exact tensor
computation below, and from a beam extension search:
$F(11)\le19$, $F(12)\le18$, $F(13)\le20$, $F(14)\le21$, $F(15)\le29$,
$F(16)\le32$, $F(17)\le32$, $F(20)\le44$.
\end{enumerate}

\section{The even-odd inequality and block floors}\label{sec:evenodd}

\begin{lemma}[Even-odd split]\label{lem:evenodd}
Let $B$ be a signing of $K_{n+m}$, let $(P,R)$ be a bipartition of the
vertices with $|P|=n$, $|R|=m$, with induced signings $A$ on $P$, $A'$
on $R$ and cross block $C\in\eps^{n\times m}$. Then for all
$x\in\eps^P$, $y\in\eps^R$,
\[
 M(B)\;\ge\;|S_A(x)+S_{A'}(y)|+|x^{\mathsf T}Cy|.
\]
\end{lemma}
\begin{proof}
$S_B(x\oplus y)=e+c$ and $S_B((-x)\oplus y)=e-c$ with
$e=S_A(x)+S_{A'}(y)$ and $c=x^{\mathsf T}Cy$, because $S_A$ is even in
$x$ while the cross term is odd. Now
$\max(|e+c|,|e-c|)=|e|+|c|$.
\end{proof}

\begin{corollary}[Block floor]\label{cor:floor}
For every signing $B$ of $K_{n+m}$ and every bipartition as above,
\[
 M(B)\;\ge\;\max_{x,y}|x^{\mathsf T}Cy|
 \;\ge\;\max\bigl(n\mu_m,\;m\mu_n\bigr),
 \qquad \mu_k:=\E\Bigl|\sum_{i=1}^k\varepsilon_i\Bigr|,
\]
where $\varepsilon_i$ are independent uniform signs. Since
$\mu_k\ge\sqrt{k/3}$ and $\mu_k=\sqrt{2k/\pi}\,(1+O(1/k))$, every
two-block construction on $n+m$ vertices pays a cross cost of order
$n\sqrt m$, and for balanced blocks a cost at least
$(\sqrt{2/\pi}-o(1))(N/2)^{3/2}=(0.28\ldots)N^{3/2}$, which is at the
target scale. Optimizing the bipartition size reproduces the July
note's bipartite-exposure bound
$\liminf G\ge\tfrac23\sqrt{2/(3\pi)}\approx0.3071$.
\end{corollary}
\begin{proof}
Take $y$ uniform and $x=\operatorname{sgn}(Cy)$ coordinatewise:
$\E_y\|Cy\|_1=\sum_{i\in P}\E|\langle c_i,y\rangle|=n\mu_m$. The bound
$\mu_k\ge k^{3/2}/\sqrt{\E S_k^4}\ge\sqrt{k/3}$ follows from
H\"older with $\E S_k^2=k$, $\E S_k^4=3k^2-2k$. The rest is
Lemma~\ref{lem:evenodd} with the even part discarded.
\end{proof}

\begin{remark}
The even part cannot currently be exploited: at the maximizing
$(x,y)$ of the cross term, the designer of $B$ can arrange the block
energies to be $O(n)$, since block energies at a fixed pair of points
have no lower bound at scale $n^{3/2}$. Converting the additive
$|e|$ term into a gain requires anticoncentration of block energies at
rounded points, which is exactly the missing landscape control
discussed in Section~\ref{sec:programs}.
\end{remark}

\section{Merging, continuity, and subsequence reductions}\label{sec:merge}

Let $q(n)$ denote the least prime $q\equiv1\pmod4$ with $q\ge n-1$. By
the prime number theorem for arithmetic progressions,
$q(n)=(1+o(1))n$.

\begin{theorem}[Merge]\label{thm:merge}
For $1\le m\le n$,
\[
 F(n+m)\;\le\;F(n)+F(m)+\sqrt{nm\,q(n)}
 \;=\;F(n)+F(m)+(1+o(1))\,n\sqrt m .
\]
\end{theorem}
\begin{proof}
Join optimal signings $A^*_n$, $A^*_m$ by a cross block $C$ equal to
any $n\times m$ submatrix of a symmetric conference matrix of order
$q(n)+1$; then
$\|C\|_{\mathrm{op}}\le\sqrt{q(n)}$ and
$|x^{\mathsf T}Cy|\le\sqrt{nm}\,\|C\|_{\mathrm{op}}$. Apply the triangle
inequality to $S=S_{A^*_n}+S_{A^*_m}+x^{\mathsf T}Cy$.
\end{proof}

By Corollary~\ref{cor:floor}, no choice of cross block can beat this
by more than the constant: the cross cost of any two-block
construction is between $(\sqrt{2/\pi}-o(1))\,n\sqrt m$ and
$(1+o(1))\,n\sqrt m$. Block composition is thus exactly a
$\Theta(n\sqrt m)$ mechanism: cheap only for $m=o(n)$, and structurally
incapable of reaching balanced scales, in agreement with the wall
recorded in the July note.

\begin{theorem}[Continuity on multiplicative windows]\label{thm:cont}
Let $\bar G=\sup_kG(k)$, which is finite. For every
$\varepsilon\in(0,1]$,
\[
 \limsup_{n\to\infty}\;
 \sup_{n\le n'\le(1+\varepsilon)n}\;|G(n')-G(n)|
 \;\le\;\Bigl(1+\tfrac52\,\bar G\Bigr)\sqrt\varepsilon .
\]
\end{theorem}
\begin{proof}
Write $n'=n+m$, $m\le\varepsilon n$. Upward:
by Theorem~\ref{thm:merge},
$G(n')\le\bigl[G(n)+G(m)(m/n)^{3/2}+(1+o(1))\sqrt{m/n}\bigr]
(1+m/n)^{-3/2}
\le G(n)+\bar G\,\varepsilon^{3/2}+(1+o(1))\sqrt\varepsilon$.
Downward: $F$ is nondecreasing (restriction of a signing does not
increase the maximum, by averaging over the last sign), so
$G(n')\ge G(n)(n/n')^{3/2}\ge G(n)(1-\tfrac32\varepsilon)$, giving
$G(n)-G(n')\le\tfrac32\bar G\,\varepsilon$. Combine, using
$\varepsilon\le\sqrt\varepsilon$.
\end{proof}

\begin{corollary}[Interval limit set]\label{cor:interval}
$|G(n+1)-G(n)|\to0$, so the set of limit points of $(G(n))$ is a
closed interval $[L^-,L^+]\subseteq[1/\pi,1/2]$. Convergence fails if
and only if $L^-<L^+$, in which case every value in $(L^-,L^+)$ is hit
infinitely often.
\end{corollary}

\begin{corollary}[Log-dense subsequences suffice]\label{cor:subseq}
Let $(n_k)$ be any sequence with $n_{k+1}/n_k\to1$, for example the
conference orders $\{q+1:q\ \text{prime},\ q\equiv1\ (\mathrm{mod}\ 4)\}$.
If $G(n_k)$ converges, then $G(n)$ converges to the same limit.
In particular the limit exists if and only if it exists along
conference orders.
\end{corollary}
\begin{proof}
Bracket $n\in[n_k,n_{k+1}]$ and apply Theorem~\ref{thm:cont} with
$\varepsilon_k=n_{k+1}/n_k-1\to0$. For the conference orders, gaps
between consecutive primes $\equiv1\pmod4$ are $o(x)$ by the prime
number theorem in arithmetic progressions.
\end{proof}

\begin{corollary}[Amplification criterion]\label{cor:amp}
Suppose there are $n_j\to\infty$ with $G(n_j)\to L:=\liminf G$ and
$\eta(n_j)\to0$ such that
\[
 F(N)\;\le\;(N/n_j)^{3/2}F(n_j)+\eta(n_j)\,N^{3/2}
 \qquad\text{for all }N\ge n_j .
\]
Then $\lim G=L$. The same holds with $\liminf$ replaced by $\limsup$
and the inequality reversed for restrictions.
\end{corollary}
\begin{proof}
$G(N)\le G(n_j)+\eta(n_j)$ for all $N\ge n_j$; let $N\to\infty$, then
$j\to\infty$.
\end{proof}

The criterion demands control from one scale to all larger scales.
Section~\ref{sec:barrier} shows that this uniformity cannot be
weakened to comparisons at any fixed ratio.

\section{The oscillation functional}\label{sec:osc}

Define, for a signing $A$ and for the minimax problem,
\[
 \operatorname{osc}(A)=\max_xS_A(x)-\min_xS_A(x),\qquad
 \Osc(n)=\min_A\operatorname{osc}(A),\qquad
 \Osc_{1/2}(n)=\tfrac12\Osc(n).
\]

\begin{proposition}\label{prop:oscbasic}
For every $A$: $\tfrac12\operatorname{osc}(A)\le M(A)\le\operatorname{osc}(A)$,
hence $F(n)/2\le\Osc_{1/2}(n)\le F(n)$.
\end{proposition}
\begin{proof}
$\sum_xS_A(x)=0$, so $\max S\ge0\ge\min S$; then
$M=\max(\max S,-\min S)$ lies between the average and the sum of the
two nonnegative numbers $\max S$ and $-\min S$. For the right-hand
inequality over minima, evaluate $\operatorname{osc}$ at an
$F$-optimal signing.
\end{proof}

\begin{theorem}[Superadditivity]\label{thm:oscsuper}
For every signing $B$ of $K_{n+m}$ and every bipartition $(P,R)$,
$\operatorname{osc}(B)\ge\operatorname{osc}(B|_P)+\operatorname{osc}(B|_R)$.
Consequently
\[
 \Osc(n+m)\ \ge\ \Osc(n)+\Osc(m)
 \qquad\text{and}\qquad
 F(n+m)\ \ge\ \Osc_{1/2}(n)+\Osc_{1/2}(m).
\]
\end{theorem}
\begin{proof}
Let $x^+$ maximize $S_{B|_P}$ and $y^+$ maximize $S_{B|_R}$. Since the
cross term is odd under $x\mapsto-x$ while both block terms are even,
one of $(\pm x^+,y^+)$ has cross term $\ge0$, so
$\max S_B\ge\max S_{B|_P}+\max S_{B|_R}$. Symmetrically
$\min S_B\le\min S_{B|_P}+\min S_{B|_R}$. Subtract. The second display
follows by evaluating at optimal signings and using
Proposition~\ref{prop:oscbasic}.
\end{proof}

\begin{theorem}[Exact value of the Gaussian-rounding method]\label{thm:varid}
For correlation matrices $R^\pm$ (unit diagonal, positive
semidefinite) and $\arcsin$ applied entrywise,
\[
 M(A)\;\ge\;\frac1{2\pi}\,
 \bigl\langle A,\ \arcsin R^+-\arcsin R^-\bigr\rangle ,
\]
and the supremum of the right side over all pairs $(R^+,R^-)$ equals
$\tfrac12\bigl(\max_xS_A-\min_xS_A\bigr)$. Hence the strongest bound
obtainable from pairs of Gaussian sign roundings is exactly
$F(n)\ge\Osc_{1/2}(n)$, and by Proposition~\ref{prop:oscbasic} this
method loses at most a factor $2$ and is exact whenever some optimal
signing has symmetric spread.
\end{theorem}
\begin{proof}
Realize $R$ as the correlation matrix of a Gaussian vector $z$, set
$x=\operatorname{sgn}(z)$, and use
$\E[\operatorname{sgn}z_i\operatorname{sgn}z_j]=\tfrac2\pi\arcsin R_{ij}$:
then $\E S_A(x)=\tfrac1\pi\langle A,\arcsin R\rangle$ up to the
$i<j$ normalization, and $|\E S_A|\le M(A)$; subtract the two choices.
For the supremum: any candidate value is
$\tfrac12(\E^+S_A-\E^-S_A)\le\tfrac12(\max S-\min S)$ since each
expectation lies in $[\min S,\max S]$; the rank-one choices
$R^\pm=x^\pm(x^\pm)^{\mathsf T}$ with $x^+$ maximizing and $x^-$
minimizing $S_A$ give
$\arcsin R^\pm=\tfrac\pi2x^\pm(x^\pm)^{\mathsf T}$ and attain it.
\end{proof}

\begin{proposition}[The linear slice and its optimality]\label{prop:slice}
The July note's choice
$R^\pm=\operatorname{corr}\bigl((I\pm tA/\sqrt n)g\bigr)$ with the
optimal $t$ yields $F(n)\ge n\sqrt{n-1}/\pi$. More generally, for any
symmetric $W$ with zero diagonal and $(W^2)_{ii}$ constant in $i$,
the slice $R^\pm=\operatorname{corr}\bigl((I\pm W)g\bigr)$ has
linear-response value
\[
 \frac2\pi\cdot
 \frac{\operatorname{tr}(AW)}
      {1+\tfrac1n\operatorname{tr}(W^2)}
 \;\le\;\frac2\pi\cdot
 \frac{\|A\|_F\,\|W\|_F}{1+\tfrac1n\|W\|_F^2}
 \;\le\;\frac{n\sqrt{n-1}}\pi ,
\]
by Cauchy--Schwarz, $\|A\|_F^2=n(n-1)$, and
$\sigma/(1+\sigma^2)\le\tfrac12$ with $\sigma^2=\|W\|_F^2/n$, with
equality only for $W\propto A$. In particular no matrix function of
$A$, and no other spectrally defined choice of $W$, improves the
constant $1/\pi$ at this order. The curvature corrections from
$\arcsin$ are nonnegative, and for quadratic slices they are
controlled by the spectral excess
$X(A)=\sum_i\lambda_i^4/(n(n-1)^2)-1\ge0$: they vanish exactly on
conference matrices, via
$\bigl\langle A,\,A B+BA\bigr\rangle=2\operatorname{tr}(A^2B)
=2n(n-1)^2X(A)$ for $B=A^2-(n-1)I$.
\end{proposition}
\begin{proof}
The first claim is the July note's Theorem 2, reverified. For the
second, expand $\Sigma^\pm=(I\pm W)^2=I\pm2W+W^2$, whose diagonal
$d=1+(W^2)_{ii}$ is constant by hypothesis, so
$\langle A,R^+-R^-\rangle=4\operatorname{tr}(AW)/d$ (the diagonal of
$A$ vanishes). When $W$ is sign-aligned with $A$
(that is, $a_{ij}W_{ij}\ge0$ for all $i<j$, as for $W=tA/\sqrt n$),
the per-edge bound $\arcsin(u+v)-\arcsin(u-v)\ge2v$ of the July note
applies with $v_{ij}=2a_{ij}W_{ij}/d\ge0$ and certifies the displayed
value; without sign alignment the per-edge bound degrades on the
misaligned edges, so the certified value is no larger. In either case
Cauchy--Schwarz bounds the display by $n\sqrt{n-1}/\pi$. The trace
identity for the curvature term is a direct computation.
\end{proof}

\begin{remark}[Why this matters]
Three consequences. First, $1/\pi$ is not a ceiling of the
Gaussian-rounding method; it is the value of its simplest slice. The
full method computes $\Osc_{1/2}(n)$, which the data below place at
$F(n)-O(1)$ for $n\le10$. Second, the adversary against every
polynomial slice is a flat-spectrum signing, and actual flat-spectrum
signings (conference matrices) have true value far above $1/\pi$
(Section~\ref{sec:attain}); the loss happens because polynomial slices
cannot see Boolean structure, only spectra. Third, replacing
$\arcsin R$ by $R$ (the semidefinite relaxation) is provably lossy:
the Grothendieck-type gap on $K_n$ is $\Theta(\log n)$
\cite{ammn,charikar-wirth}, so any argument through the linearized
functional must lose a $\log n$ factor in general; the arcsine
nonlinearity, or the restriction to the minimizing $A$, must be used.
\end{remark}

Enumeration gives the oscillation optima
\[
\begin{array}{c|cccccccc}
n&3&4&5&6&7&8&9&10\\\hline
\Osc(n)&4&8&8&10&16&20&24&26\\
2F(n)&6&8&8&10&18&20&24&26
\end{array}
\]
so $\Osc=2F$ whenever parity permits and $\Osc=2F-2$ otherwise
($n=3,7$). The superadditivity of Theorem~\ref{thm:oscsuper} is
strict in the data.

\begin{conjecture}[Symmetric spread]\label{conj:spread}
$\Osc_{1/2}(n)=F(n)-O(n)$; equivalently, some near-optimal signings
have $\max S\approx-\min S$. Under this conjecture $F$ is
asymptotically superadditive:
$F(n+m)\ge F(n)+F(m)-O(n)$.
\end{conjecture}

Even granting Conjecture~\ref{conj:spread}, superadditivity does not
force convergence at the superextensive scale $n^{3/2}$: summing
$k$ blocks of size $n$ captures only $k\,n^{3/2}$ of the
$(kn)^{3/2}=k^{3/2}n^{3/2}$ target, so Fekete-type arguments are
vacuous in both directions. This is quantified in
Section~\ref{sec:barrier}.

\section{Fixed-temperature equivalence}\label{sec:beta}

\begin{proposition}\label{prop:beta}
For $\beta>0$ let
$\Phi_{n,\beta}
=\min_A\frac1{\beta n^{3/2}}\log\sum_{x\in\eps^n}e^{\beta|S_A(x)|}$.
Then
\[
 0\;\le\;\Phi_{n,\beta}-G(n)\;\le\;\frac{\log2}{\beta\sqrt n}.
\]
Hence for every fixed $\beta>0$ (indeed whenever
$\beta_n\sqrt n\to\infty$), $\liminf$, $\limsup$, and existence of the
limit coincide for $\Phi_{n,\beta}$ and $G(n)$. The same holds for the
sign-augmented partition function
$\sum_{x,s}e^{\beta sS_A(x)}$ with error $\log(2^{n+1})/(\beta n^{3/2})$.
\end{proposition}
\begin{proof}
$e^{\beta M}\le\sum_xe^{\beta|S|}\le2^ne^{\beta M}$ for each $A$; take
$\log$, divide, minimize.
\end{proof}

Two of the concerns listed against the finite-temperature route are
therefore empty: no uniformity in $\beta$ is needed, and the absolute
value can be handled exactly. What remains open is precisely the
$n\to\infty$ behavior at fixed $\beta$ with adversarial couplings. The
July note's annealed obstruction is confirmed quantitatively: for a
balanced merge the annealed cross cost is
$nm\log\cosh\beta/(\beta N^{3/2})\asymp\beta N^{1/2}$, which requires
$\beta\sqrt N\to0$, while the sandwich above requires
$\beta\sqrt N\to\infty$. The window is empty, and this is structural
for annealing, not an artifact.

\section{Relaxations collapse}\label{sec:relax}

\begin{proposition}[Box and sphere]\label{prop:relax}
Let $m=\binom n2$ and let $\chi(x)=(x_ix_j)_{i<j}$.
\begin{enumerate}
\item $\displaystyle\min_{a\in[-1,1]^m}\max_x|\langle a,\chi(x)\rangle|=0$,
attained at $a=0$; every convex relaxation of the coupling cube whose
feasible set contains $0$ has optimum $0$.
\item $\displaystyle\min_{\|a\|_2=\sqrt m}\max_x|\langle a,\chi(x)\rangle|
=\sqrt m$, attained by concentrating all mass on a single edge:
$|\langle a,\chi\rangle|\equiv\sqrt m$ for $a=\sqrt m\,e_{\{1,2\}}$.
The lower bound is
$\max_x\langle a,\chi\rangle^2\ge\E_x\langle a,\chi\rangle^2=\|a\|_2^2$.
\end{enumerate}
\end{proposition}

So the $\pm1$ coupling constraint plays two distinct roles at once:
it fixes the total mass and it forces the mass to spread over all
$m$ edges. Dropping either role collapses the value from
$\Theta(n^{3/2})$ to $O(n)$. The integrality gap of any relaxation of
this kind is the entire phenomenon, and rounding a relaxed optimizer
back to signs must inject, not preserve, the $n^{3/2}$ scale. This
disposes of the relaxed-couplings attack in its natural formulations.

\begin{theorem}[Zero graphon]\label{thm:graphon}
If $M(A_n)=O(n^{3/2})$, then the cut norm satisfies
$\|A_n\|_\square\le4M(A_n)=O(n^{3/2})=o(n^2)$: every sequence of
near-optimal signings converges to the zero graphon, and every
functional continuous in the dense graph (cut-metric) topology is
blind to the problem. Any limit theory that could express
$\lim G$ must live at the $n^{3/2}$ fluctuation scale, for which no
compactness theorem is currently available.
\end{theorem}
\begin{proof}
$\|A\|_\square\le\max_{x,y\in\eps^n}|x^{\mathsf T}Ay|$ by writing
indicators as $(1+x)/2$; polarization and multilinearity of the
zero-diagonal quadratic form give
$\max_{x,y}|x^{\mathsf T}Ay|\le2\max_z|z^{\mathsf T}Az|=4M(A)$.
\end{proof}

\section{Conference matrices at their ceiling}\label{sec:attain}

\begin{theorem}[Attainment]\label{thm:attain}
Let $C$ be a symmetric conference matrix of order $N$. Then
$M(C)\le\tfrac12N\sqrt{N-1}$, with equality if and only if some
$x\in\eps^N$ is an eigenvector of $C$, necessarily with eigenvalue
$\pm\sqrt{N-1}\in\mathbb Z$. Equality holds at $N=10$ (Petersen
two-graph, $x=\mathbf1$ for the Petersen-form Seidel matrix, value
$15$) and at $N=26$ (Paley over $\mathbb F_{25}$, value $65$), both
verified exactly. Equality is impossible whenever $N-1$ is not a
perfect square, in particular at the orders $6,14,18,30$ checked in
the July note, where the exact values are $5,21,33,75$ and the
normalized ratios $0.340,0.401,0.432,0.456$ increase toward $1/2$.
\end{theorem}
\begin{proof}
Cauchy--Schwarz: $|x^{\mathsf T}Cx|\le\|x\|\,\|Cx\|=N\sqrt{N-1}$, with
equality iff $Cx\parallel x$; integrality of $x^{\mathsf T}Cx$ forces
$\sqrt{N-1}\in\mathbb Q$. A $\pm1$ eigenvector with eigenvalue
$\sqrt{N-1}$ exists iff the switching class contains a graph that is
regular of degree $(N-1-\sqrt{N-1})/2$ (switch by the eigenvector and
read off constant row sums): degree $3$ at $N=10$ (Petersen), degree
$10$ at $N=26$ (a Paulus graph). The stated values are exhaustive
computations with recorded witnesses.
\end{proof}

Consequences.

\begin{enumerate}
\item There is no uniform conference deficit. The July note's wall
``small Paley examples do not attain the spectral ceiling'' is
explained entirely by $\sqrt{N-1}\notin\mathbb Z$ at those orders. At
orders $q+1$ with $q$ an even power of an odd prime, whenever the
switching class contains a regular graph of degree $(q-\sqrt q)/2$,
the construction value is exactly $\tfrac12(q+1)\sqrt q$, of
normalized ratio $\tfrac12\sqrt{q/(q+1)}\to\tfrac12$. Candidate orders
$10,26,50,82,\dots$; whether infinitely many contain such regular
descendants is a known open-ended question in the theory of regular
two-graphs.
\item Improving $\limsup G\le\tfrac12$ therefore requires either
non-conference matrices or a proof that the minimum beats the
conference class at scale. The datapoint $F(10)=13<15$ shows the
minimum does beat the conference class at $n=10$, by a non-flat
signing. The upper-bound program should not fixate on spectrally flat
matrices.
\item The certification barrier remains: for a random signing the true
maximum is $(0.7632\ldots)n^{3/2}$ (the Parisi ground state, universal
for $\pm1$ couplings \cite{parisi-gs,gt}), while no known efficient
certificate goes below the spectral $n^{3/2}$; and for structured
matrices every known certificate below the spectral value is exhausted
by exact enumeration. New certification ideas (Krawtchouk or Delsarte
arguments on the coset structure of the extended cut code) are listed
as an open direction, with the caution that the code's dual weight
spectrum is too rich for the standard external-distance bound.
\end{enumerate}

\section{The consistency barrier}\label{sec:barrier}

Everything proved so far constrains $G$ only locally in scale. The
next proposition shows this is a hard limitation of the entire
inequality set, not of the sharpness of individual constants.

\begin{proposition}[A divergent profile consistent with all proved
inequalities]\label{prop:barrier}
Define
\[
 \widetilde G(n)=\frac25+\frac1{20}
 \sin\bigl(2\pi\log_2\log_2n\bigr)\quad(n\ge4),
 \qquad
 \widetilde F(n)=\widetilde G(n)\,n^{3/2}.
\]
Then, for all sufficiently large $n$ and all $1\le m\le n$:
$\widetilde F$ is strictly increasing;
$\widetilde F(n+1)\le\widetilde F(n)+n$;
$\widetilde F(n+m)\le\widetilde F(n)+\widetilde F(m)+1.01\,n\sqrt m$;
$\widetilde F(n+m)\ge\widetilde F(n)+\widetilde F(m)$;
$\widetilde F(n)/n^{3/2}\in[1/\pi,1/2]$;
and $\widetilde G$ satisfies the conclusions of
Theorems~\ref{thm:cont} and Corollaries~\ref{cor:interval},
\ref{cor:subseq} vacuously or with margin. Yet
$\widetilde G(n)$ diverges: its limit set is $[0.35,0.45]$.
\end{proposition}
\begin{proof}
All verifications are elementary calculus from
$\widetilde G'(n)=O\bigl(1/(n\log n)\bigr)$, so that
$\widetilde G$ is asymptotically constant on any window of bounded
ratio, together with
$\tfrac d{dn}\widetilde F=n^{1/2}\bigl(\tfrac32\widetilde G+o(1)\bigr)
\in[0.5\,n^{1/2},\,0.7\,n^{1/2}]$. Monotonicity and the Lipschitz
bound follow at once. For the merge inequality, the left side minus
$\widetilde F(n)$ is $\tfrac32\widetilde G\,\varepsilon n^{3/2}(1+o(1))$
with $\varepsilon=m/n$, which is at most
$\widetilde G(m)m^{3/2}+1.01\,n\sqrt m$ because
$0.7\varepsilon\le1.01\sqrt\varepsilon$; superadditivity holds because
$0.5\varepsilon n^{3/2}\ge0.45\,\varepsilon^{3/2}n^{3/2}$ for
$\varepsilon\le1$. The oscillation of $\log_2\log_2n$ is unbounded
while its increment over $[n,2n]$ is $O(1/\log n)$, so the limit set
is the full interval. Numerical verification of every inequality with
explicit margins for $n\le10^6$ is included in the verification
suite.
\end{proof}

The profile also satisfies the oscillation chain of
Section~\ref{sec:osc} with $\widetilde\Osc:=2\widetilde F$, and
Conjecture~\ref{conj:spread}. Conclusion: the union of all inequalities
proved in this repository, even sharpened by constants, cannot decide
the question. Deciding it requires an inequality that couples scale
$n$ to scales $N$ with $N/n\to\infty$, such as the amplification
criterion of Corollary~\ref{cor:amp}. Comparisons at any fixed ratio
$\lambda$, for instance a proof that $G(2n)-G(n)\to0$, would still not
suffice: $\widetilde G$ satisfies that too.

\section{Three proof programs}\label{sec:programs}

\subsection*{Program 1: elliptope slices (lower bound)}
Target lemma: a map $A\mapsto(R^+(A),R^-(A))$ into correlation-matrix
pairs, analyzable uniformly in $A$, with
$\tfrac1{2\pi}\langle A,\arcsin R^+-\arcsin R^-\rangle
\ge(1/\pi+\delta)\,n^{3/2}$ for some $\delta>0$. By
Theorem~\ref{thm:varid} the program is complete in principle: the full
supremum equals $\Osc_{1/2}(n)\ge F(n)/2$, and empirically
$\Osc_{1/2}=F-O(1)$. Failure points identified: every polynomial slice
degenerates on flat-spectrum signings
(Proposition~\ref{prop:slice}); rank-one slices require knowing
extremal Boolean vectors, which is circular; the linearized functional
is off by $\Theta(\log n)$ in general \cite{ammn}. Repair directions:
mixtures of a polynomial slice with low-rank Boolean-structured
corrections, controlled by an anticoncentration statement for block
energies at rounded points (the missing landscape lemma), or slices
adapted to the bipartite exposure structure of
Lemma~\ref{lem:evenodd} (Rademacher rather than Gaussian roundings).
The obstruction appears structural for slices defined by spectral data
alone, and technical, possibly movable, for Boolean-aware slices.

\subsection*{Program 2: amplification (upper bound)}
Target lemma (Corollary~\ref{cor:amp}): from one good signing at size
$n$, signings at all $N\ge\Lambda(n)$ with
$G(N)\le G(n)+o_n(1)$. Exact losses of the known lifts:
block merges cost $\Theta(n\sqrt m)$ exactly
(Theorem~\ref{thm:merge} and Corollary~\ref{cor:floor}), fatal at
balanced scales; iterated one-vertex extension is far weaker in
practice (beam search from the exact $F(10)$ optimizer reaches only
ratio $\approx0.52$ by $n=14$, worse than conference restrictions);
tensor doubling $A\otimes H_2$ (intra-pair edges refilled) does not
preserve the constant: exact values move the ratios
$0.358\to0.411$, $0.340\to0.433$, $0.486\to0.477$, $0.442\to0.500$,
$0.411\to0.492$ at $n=5,6,7,8,10$. The data are consistent with
iterated tensoring converging to its own fixed constant near $1/2$
rather than preserving the base ratio: misaligned block vectors let
the adversary realize Gram-relaxation values, and the loss is a
Grothendieck-type constant, not $o(1)$. Repair directions: lifts along
structured algebraic extensions (subfield towers of Paley matrices,
where the added symmetry may rule out misaligned adversaries), or
merge trees whose cross blocks are chosen adaptively against the
blocks' own near-maximizers rather than obliviously. The oblivious
version of the obstruction is proved (Corollary~\ref{cor:floor});
whether adaptive cross blocks can beat $n\sqrt m$ is open and is the
sharpest open sub-question of this program.

\subsection*{Program 3: adversarial thermodynamic interpolation}
Target: at fixed $\beta$, an interpolation between $\Phi_{N,\beta}$
and block free energies with normalized cost $o(1)$. Status: the
sandwich of Proposition~\ref{prop:beta} makes the temperature exchange
exact; the annealed route has an empty parameter window (verified
quantitatively in Section~\ref{sec:beta}); the Guerra--Toninelli
mechanism \cite{gt} needs a disorder average and Gaussian integration
by parts to produce a signed overlap term, and the minimum over
couplings supplies neither: the derivative of a minimum runs through
the argmin, whose structure is exactly what is unknown. A quenched
interpolation with adversarial couplings would need a new convexity
input in the coupling variable. Note that the natural convexity is
absent: $A\mapsto\log\sum e^{\beta|S_A|}$ is convex, so its minimum
over the cube is attained at vertices (consistent with the collapse of
Proposition~\ref{prop:relax}), but minima of convex functions do not
interpolate. The obstruction appears structural.

\section{The sharpest wall and the next lemma}\label{sec:wall}

The problem is now reduced, with proofs, to the following statement.
All local behavior of $G$ is understood at the level needed: $G$ has
an interval limit set inside $[1/\pi,1/2]$, is continuous on bounded
multiplicative windows, and its limit exists iff it exists along
conference orders. Every additive composition mechanism is exactly
priced: merges cost $\Theta(n\sqrt m)$, oscillation is superadditive,
and the superextensive scale makes both Fekete directions vacuous. The
value class of all proved inequalities admits divergent profiles
(Proposition~\ref{prop:barrier}), so no refinement of these local
estimates decides the question.

What is missing is a single scale-coupling inequality. The recommended
target, the most constrained version consistent with everything above:

\begin{conjecture}[Amplification lemma]\label{conj:amp}
There is a function $\eta$ with $\eta(n)\to0$ such that for every $n$
and every $N\ge n^2$,
\[
 F(N)\;\le\;(N/n)^{3/2}\,F(n)\;+\;\eta(n)\,N^{3/2}.
\]
\end{conjecture}

By Corollary~\ref{cor:amp} this implies convergence. The exponent
$2$ in $N\ge n^2$ is inessential; any threshold $\Lambda(n)<\infty$
works. The evidence: tensor and merge lifts fail by constants, not by
scale, and the failure constant shrinks as the base ratio approaches
the tensor fixed point; the conjecture only requires lossless
amplification from optimal bases, whose structure (exact witnesses
above) is rigid and highly symmetric at small $n$. The first concrete
test: construct, from the exact $F(10)$ optimizer or the order-$26$
conference class, signings at $n\approx10^2$ with ratio provably below
$0.46$. A disproof of Conjecture~\ref{conj:amp} in a strong form
(a lower-bound mechanism active only at special scales) is the only
currently visible route to nonconvergence, and nothing in the data
suggests such a mechanism: the upper envelope moves smoothly, and all
known algebraic families exist at log-dense orders, which
Corollary~\ref{cor:subseq} then smooths out.

\appendix

\section{Computational protocol}

All computations are exhaustive or witness-carrying; none are
heuristic except the beam search, which only produces upper bounds
that are then verified exactly.

\begin{itemize}
\item \texttt{exact\_fn.c}: exact $F(n)$, one-sided optimum, and
$\Osc(n)$ for $n\le10$ by Gray-code enumeration of all switching
classes (normalized representatives), with periodic from-scratch
recomputation and a final independent audit of reported minima.
$n=10$ (about $6.9\cdot10^{10}$ class evaluations) was run as $8$
disjoint fixed-prefix slices, each auditing its own minima.
\item \texttt{conference\_max.c}: symmetric conference matrices at
orders $6,10,14,18,26,30$ (Paley over prime fields, Petersen Seidel
form, Paley over $\mathbb F_{25}$), verified to satisfy
$C^2=(N-1)I$, exact Boolean maxima with witnesses, signed spreads,
and nested restriction chains inside order $30$.
\item \texttt{ext\_chain.c}: beam extension upper bounds from the
exact $F(10)$ witness.
\item \texttt{verify\_new\_results.py}: independent reverification of
every table in this note from recorded witnesses, full independent
enumeration for $n\le6$, parity checks, Petersdorf check, tensor
recomputation, the finite merge instance, numerical margins for
Proposition~\ref{prop:barrier} up to $n=10^6$, and corruption
controls that must fail when data are perturbed.
\end{itemize}

Restriction-chain data (order 30, three random nested chains):
normalized ratios decrease monotonically in the subgraph size from
about $0.7$ at $k=8$ to $0.456$ at $k=30$. Restrictions of good large
signings are relatively poor small signings; the naive
restriction-decay route to the reversed amplification criterion fails
on conference matrices.

\begin{thebibliography}{9}

\bibitem{petersdorf}
M.~Petersdorf,
\emph{Einige Bemerkungen \"uber vollst\"andige Bigraphen},
Wiss. Z. Techn. Hochsch. Ilmenau \textbf{12} (1966), 257--260.

\bibitem{ammn}
N.~Alon, K.~Makarychev, Y.~Makarychev, and A.~Naor,
\emph{Quadratic forms on graphs},
Invent. Math. \textbf{163} (2006), 499--522.

\bibitem{charikar-wirth}
M.~Charikar and A.~Wirth,
\emph{Maximizing quadratic programs: extending Grothendieck's
inequality}, FOCS 2004, 54--60.

\bibitem{parisi-gs}
A.~Auffinger and W.-K.~Chen,
\emph{Parisi formula for the ground state energy in the mixed p-spin
model}, Ann. Probab. \textbf{45} (2017), 4617--4631.

\bibitem{gt}
F.~Guerra and F.~L.~Toninelli,
\emph{The thermodynamic limit in mean field spin glass models},
Comm. Math. Phys. \textbf{230} (2002), 71--79.

\bibitem{sole-zaslavsky}
P.~Sol\'e and T.~Zaslavsky,
\emph{A coding approach to signed graphs},
SIAM J. Discrete Math. \textbf{7} (1994), 544--553.

\bibitem{dmp}
A.~Defant, M.~Masty\l{}o, and A.~P\'erez,
\emph{On the Fourier spectrum of functions on Boolean cubes},
Math. Ann. \textbf{374} (2019), 653--680.

\bibitem{seidel}
J.~J.~Seidel,
\emph{A survey of two-graphs},
Colloquio Internazionale sulle Teorie Combinatorie (Rome, 1973),
Tomo I, 481--511.

\end{thebibliography}

\end{document}
